\section{Axiom y los Neo Lab: qué requiere emprender en AI for Math}

La descripción del video menciona que Axiom cerró una ronda Serie A de 200 millones de dólares con una valuación de 1,600 millones de dólares; también llamó la atención que Ken Ono dejara o pusiera en pausa la trayectoria tradicional de profesor para unirse a Axiom. Más allá de cómo lo narren los medios, esto muestra que AI for Math pasó de ser un problema académico a ser un problema de Neo Lab con alta densidad de capital.

\begin{figure}[H]
\centering
\includegraphics[width=0.92\textwidth]{ai-math-neolab.png}
\caption{Componentes de un Neo Lab de AI for Math: matemáticos, sistemas de formalización, entrenamiento de modelos, financiamiento y capacidad de cómputo, y objetivos de producto.}
\end{figure}

\begin{knowledgebox}{Lectura de la figura: un Neo Lab no es una empresa emergente común}
Los cinco elementos de la figura apuntan en conjunto a organizaciones como Axiom: la red de matemáticos de primer nivel aporta los problemas y el criterio estético, Lean/Coq aporta la base de formalización, el entrenamiento de modelos aporta la capacidad de automatización, el financiamiento y la capacidad de cómputo sostienen los experimentos a escala, y los objetivos de producto convierten la investigación en sistemas utilizables.
\end{knowledgebox}

\subsection{Por qué importa la red de matemáticos}

Los datos matemáticos no son datos web comunes. Los problemas realmente valiosos, las rutas de demostración, la organización de las bibliotecas de teoremas y el juicio estético suelen estar en manos de la comunidad de matemáticos. Que se sume un matemático como Ken Ono no es solo un respaldo: es probable que aporte la selección de problemas, una red de investigación y criterios de verificación.

\begin{importantbox}{El círculo virtuoso de datos de AI for Math}
Si el sistema puede ayudar a los matemáticos a formalizar demostraciones, completar lemmas y encontrar contraejemplos, y los matemáticos a su vez pueden devolver correcciones y evaluaciones, podría formarse un círculo virtuoso de datos. Sin la participación de matemáticos, el modelo tiende a quedarse en bancos de problemas y demostraciones superficiales.
\end{importantbox}

\subsection{La emprendedora menos probable}

En la entrevista, Hong Letong (洪乐潼) habla de «la emprendedora menos probable». Este tipo de Neo Lab es distinto del emprendimiento tradicional de aplicaciones: quien lo funda debe entender matemáticas, IA, sistemas de formalización, organización de la investigación y financiamiento. Se parece más a comprimir un instituto de investigación, un equipo de modelos y una empresa de producto en una sola organización de alta densidad.

\begin{warningbox}{El doble riesgo del emprendimiento orientado a la investigación}
Si es demasiado académico, la ruta hacia el producto y los ingresos no queda clara; si se orienta demasiado al producto, puede perder profundidad en la investigación básica. AI for Math necesita mantener la tensión entre los objetivos científicos y la implementación de ingeniería.
\end{warningbox}

\subsection{Resumen de la sección}

Axiom representa un nuevo tipo de AI Lab: los problemas de investigación son lo bastante fundamentales, la densidad de capital es lo bastante alta y el equipo necesita conectar a la vez a los matemáticos, los sistemas de formalización y el entrenamiento de modelos. No es un emprendimiento SaaS común ni un laboratorio académico tradicional.
